\subsection{Stage~K (Cross-Node Alternative): kind Multi-Worker CNI Sensitivity}
\label{sec:analysis-rq14b}

The Stage~K cross-node alternative is an explicitly supporting sensitivity stage rather than a
primary thesis-facing measurement: it re-runs the Stage~K chain family
on a six-worker \texttt{kind} cluster with pod anti-affinity forcing
every chain segment onto a distinct kind worker container and the
benchmark client onto a sixth dedicated worker
(\cref{sec:methods-rq14b}). The analytical contribution is bounded
by a methodological reading carried over from
\cref{sec:methods-rq14b}: kind's worker ``nodes'' are Docker
containers that share one host kernel and communicate over kind's
CNI bridge; they are not separate physical or virtual hosts. Stage~K (cross-node)
therefore measures the cost of routing inter-service gRPC through
the kind CNI bridge under anti-affinity, not the cost of real
cross-host networking. The genuine multi-host inter-node penalty for
the same chain family is measured in
\cref{sec:analysis-rq15b}; the role of Stage~K (cross-node) is to bound how much
the single-node placement of Stage~K hides on the \emph{same physical
host}, and to test whether the Stage~K service-count ordering survives
when the chain traffic is forced through the kind CNI bridge.

\paragraph{The service-count ordering and the bounded-non-linearity pattern survive the kind-CNI probe.}
Under the kind multi-worker condition, mean end-to-end latency
increases monotonically from 67.491~ms for
\texttt{monolithic\_k8s\_1svc} to 87.378~ms, 120.465~ms, 135.639~ms,
and 141.169~ms for \texttt{chain\_2svc} through \texttt{chain\_5svc}
(\cref{tab:rq14b_results}). The same monotonic ordering established
in Stage~K is therefore preserved when each chain hop is forced across
a kind CNI-bridge boundary on the same host. The
bounded-non-linearity reading of Stage~K is also preserved: the final
\texttt{chain\_4svc}~$\rightarrow$~\texttt{chain\_5svc} step adds
approximately 5.530~ms, smaller than the preceding +15.174~ms
\texttt{chain\_3svc}~$\rightarrow$~\texttt{chain\_4svc} step, in line
with the smaller 98~KiB \texttt{layer4} activation added at the final
boundary (the runner's \texttt{total\_activation\_bytes\_mean} rises
from 1960~KiB to 2058~KiB across these two conditions, matching the
Stage~K profile). The reading is therefore narrow but useful: the
Stage~K chain-cost gradient and the activation-driven shape of the
marginal cost are not specific to single-node pod colocation; they
persist once chain traffic is routed through kind's CNI bridge.

\paragraph{The kind multi-worker per-condition penalty is an order of magnitude larger than the AKS multi-node penalty, confirming that kind multi-node is not a substitute for cross-host networking.}
Computed condition-by-condition against the frozen Stage~K single-node
baseline, the per-condition penalty
$\Delta_c = \mathrm{mean}_c^{\mathrm{kind\text{-}cni}} - \mathrm{mean}_c^{\mathrm{single}}$
is +3.991~ms for \texttt{chain\_2svc}, +32.737~ms for
\texttt{chain\_3svc}, +45.067~ms for \texttt{chain\_4svc}, and
+48.634~ms for \texttt{chain\_5svc} (\cref{tab:rq14b_delta}). The
monolithic condition is faster on average in Stage~K (cross-node) ($-13.575$~ms)
but its per-iteration standard deviation rises from 4.017~ms (Stage~K)
to 10.999~ms (\cref{tab:rq14b_results}), its median (61.917~ms) is
well below its mean (67.491~ms), and its per-round means are
$\{62.000, 68.553, 68.677, 69.107, 69.119\}$~ms (\cref{tab:rq14b_rounds},
the first round is markedly faster than rounds 2--5). This
monolithic shift is therefore consistent with a one-off cluster-start
artefact and host co-tenancy variability, and should not be read as
evidence that kind multi-worker placement is faster than single-node
placement. For the chained conditions, the kind-CNI per-condition
deltas are roughly an order of magnitude larger than the
corresponding multi-node AKS deltas reported in
\cref{sec:analysis-rq15b} (+2.131~ms to +4.086~ms across the chained
conditions on AKS, \cref{tab:rq15b_delta}). This separation is the
key analytical contribution of Stage~K (cross-node): kind's ``multi-node''
topology is not a substitute for genuine cross-host networking. CNI
plugin choice and the container-network path are an independent
source of measured-latency variance in Kubernetes
clusters~\cite{cni_k8s_ic2e21}, and absolute timings are known not
to transfer cleanly across container, VM, and bare-metal execution
substrates~\cite{virtualization_costs}; treating the kind-multi
result as a cross-host estimate would inflate the inferred
inter-node cost relative to what \cref{sec:analysis-rq15b} actually
measures on AKS.

\paragraph{Round-level dispersion inflates relative to Stage~K and is not chain-explained.}
Per-condition round CVs in Stage~K (cross-node) (\cref{tab:rq14b_rounds}) are
0.0290 (\texttt{chain\_2svc}), 0.0551 (\texttt{chain\_3svc}), 0.0187
(\texttt{chain\_4svc}), 0.0247 (\texttt{chain\_5svc}), and 0.0456
(\texttt{monolithic\_k8s\_1svc}), compared with the much tighter
0.0074--0.0161 range observed in single-node Stage~K
(\cref{tab:rq14_rounds}). Two features matter for interpretation.
First, the highest CVs do not track chain depth: monolithic (no
inter-pod traffic) and \texttt{chain\_3svc} (three services) are
noisier than \texttt{chain\_5svc} (five services), so the inflated
CV is not a property of having more chain hops. Second, the
monolithic round-mean trajectory above is consistent with a
first-round warm-up shift, not with a steady-state difference. Both
observations point to kind-CNI variability and shared-host effects
on a single physical machine that hosts all six kind worker
containers, rather than a chain-specific noise mechanism; this
reading is consistent with the broader observation that orchestrator-
and host-level contention can dominate measured deltas on shared
infrastructure~\cite{noisy_neighbor_k8s}. By the SC'15 reporting
standard, round-level consistency is the primary indicator of result
stability, and the round CVs above are too high to read Stage~K (cross-node) as a
primary number~\cite{benchmarking_sc15}: the within-RQ ordering and
the qualitative bounded-non-linearity claim remain interpretable,
but the absolute millisecond magnitudes should be quoted only as
upper sensitivity bounds, not as deployment-grade values.

\paragraph{Validation.}
Three pieces of evidence support a within-run reading of Stage~K (cross-node).
First, functional parity holds: both the local segment-chain check
and the live in-cluster gRPC check report
\texttt{max\_abs\_diff}~$=0.0$ at the predeclared tolerance of
$1.0\times10^{-5}$ across all five conditions
(\texttt{merged\_results/parity\_validation.json}, summarised
alongside \cref{tab:rq14b_results}), so the within-Stage~K (cross-node) ordering
is not confounded by numerical drift across the partitioned graph.
Second, infrastructure controls are applied symmetrically across
conditions: image (\texttt{thesis-inference:latest}, pull policy
\texttt{IfNotPresent}), namespace (\texttt{rq14b}), resource
requests and limits (1~CPU, 512~MiB per service pod; 1~CPU, 1~GiB
client), threading (one thread per pool), and CPU affinity to
core~0 are uniform across the five conditions
(\texttt{merged\_results/deployment\_metadata.json},
\texttt{merged\_results/environment.json}). Third, multi-node
placement was enforced by the same anti-affinity rule used in
Stage~C (multi-node), with six distinct worker containers and a dedicated client
worker. Four limits must be stated explicitly. (i) The image tag
\texttt{thesis-inference:latest} is a mutable Docker tag, so
byte-exact reproduction is bounded by whatever image was loaded
into the kind cluster at run time, the same caveat that applies to
Stage~K (\cref{sec:methods-artifacts}). (ii) Host CPU controls were
weaker than in Stage~K: governor and turbo writes to \texttt{/sys}
failed because those paths were read-only in the containerised
environment, and the detected governor was \texttt{schedutil} with
turbo enabled (rather than the \texttt{performance} governor and
disabled turbo observed for Stage~K). The \texttt{nice~-5} request
also failed for permission reasons. Stage~K (cross-node) is therefore not a clean
replication of Stage~K with only a topology change; reduced host
control is one plausible co-contributor to the higher round CVs.
(iii) Comparisons against Stage~K are interpreted at the
\emph{per-condition} delta level (\cref{tab:rq14b_delta}), not as a
cross-stage means comparison, because the two stages differ in host
control state in addition to placement. (iv) Stage~K (cross-node) is one host,
one kind cluster, one fixed seed; the kind worker containers share
a kernel and the same physical hardware, so cross-host replication
is, by construction, out of scope and is addressed only by Stage~C (multi-node)
on AKS.

\paragraph{Evaluation.}
The evaluative reading of Stage~K (cross-node) has three facets. First, the
result is sufficient to support a qualitative claim: the Stage~K
chain-cost ordering and the activation-driven bounded-non-linearity
pattern are not specific to single-node pod colocation; they
survive once every chain hop is forced through the kind CNI bridge.
Second, the result is sufficient to bound the additional cost a
switch to a kind multi-worker topology would add to the single-node
Stage~K measurement on this host: between +3.991~ms
(\texttt{chain\_2svc}) and +48.634~ms (\texttt{chain\_5svc}) per
chained condition (\cref{tab:rq14b_delta}), with the worst case at
the deepest chain.
Third, the result is not sufficient to support a deployment-grade
multi-host claim, a production-readiness claim, or a scalability
claim: the per-condition deltas are an order of magnitude larger
than the AKS-measured multi-node deltas of
\cref{sec:analysis-rq15b}, the host-control state was degraded
relative to Stage~K, and the round CVs are correspondingly inflated.
The headline Stage~K evaluative conclusion is unchanged; Stage~K (cross-node) adds
a bounded sensitivity caveat, not an independent deployment claim,
and motivates Stage~C (multi-node) as the appropriate evidence path for genuine
inter-node cost.

\paragraph{Stage~K (cross-node) outcome.}
Under the kind multi-worker CNI-bridge probe, the Stage~K
architectural ordering is preserved: mean end-to-end latency
increases monotonically from 67.491~ms for
\texttt{monolithic\_k8s\_1svc} through 87.378~ms, 120.465~ms,
135.639~ms, and 141.169~ms for \texttt{chain\_2svc} through
\texttt{chain\_5svc} (\cref{tab:rq14b_results}), and the
bounded-non-linearity pattern survives (the
\texttt{chain\_4svc}~$\rightarrow$~\texttt{chain\_5svc} increment of
+5.530~ms is smaller than the preceding +15.174~ms
\texttt{chain\_3svc}~$\rightarrow$~\texttt{chain\_4svc} step,
matching the small 98~KiB \texttt{layer4} activation moved at the
added boundary). The per-condition kind-CNI penalty relative to the
single-node Stage~K baseline ranges from +3.991~ms to +48.634~ms
across the chained conditions (\cref{tab:rq14b_delta}), an order of
magnitude larger than the AKS-measured multi-node penalty reported
in \cref{sec:analysis-rq15b}, and round-level CVs (up to 0.0551 for
\texttt{chain\_3svc} and 0.0456 for the monolithic baseline,
\cref{tab:rq14b_rounds}) are materially inflated relative to Stage~K.
The contribution of Stage~K (cross-node) to the deployment analysis is therefore
two-fold and bounded: (i) it shows that the Stage~K ordering and the
activation-driven bounded-non-linearity are not artefacts of
single-node colocation on the same host, and (ii) it shows that
kind multi-node deltas are not a substitute for a managed-cloud
cross-host measurement and should be read as a kind-CNI sensitivity
envelope rather than as a deployment-grade multi-node result.
